\section{Kraft 不等式与 Huffman 编码}


上一节给出了前缀码平均码长的上下界；这一节说明“哪些码长组合可以实现”，并给出达到最优的构造算法：\textbf{Huffman 编码}。

\textbf{本节用到的记号：} 码长 $l_1,\dots,l_m$；Kraft 和 $K=\sum_i2^{-l_i}$；$H(P)$ 为熵；$L^*$ 为最优前缀码的平均码长。

\subsection{Kraft 不等式}


\textbf{定理.} 存在以 $l_1,\dots,l_m$ 为码长的前缀码，当且仅当

\[
\boxed{
\sum_{i=1}^m2^{-l_i}\le1
}.
\]

\textbf{必要性证明（随机无限串）.} 独立均匀地生成一个无限比特串 $S\in\{0,1\}^{\mathbb N}$。对每个码字 $c_i$，事件“$c_i$ 是 $S$ 的前缀”的概率是

\[
2^{-l_i}.
\]

前缀码保证这些事件两两不交（一个无限串不可能同时以两个不同码字为前缀）。因此

\[
\sum_{i=1}^m2^{-l_i}
=\sum_i\Pr[c_i\text{ 是 }S\text{ 的前缀}]
\le1.
\]

\textbf{充分性（构造树）.} 把长度 $l_i$ 的码字看成深度为 $l_i$ 的叶子。按长度从小到大放置：每个长度为 $l_i$ 的叶子占据一个“容量为 $2^{-l_i}$”的区间；Kraft 和 $\le1$ 保证这些区间不相交，于是可以构造出一棵前缀码树。

\textbf{直观.} Kraft 和就是所有码字叶子在无限比特串空间中所占的“总测度”；前缀码的叶子不能重叠，所以总测度不超过 $1$。

\subsection{Huffman 编码}


\textbf{算法.}

1. 把每个符号 $x_i$ 看成一个权为 $P(x_i)$ 的叶子；
2. 每次选出权最小的两个结点，合并成一个新结点，权为两者之和；
3. 重复直到只剩一个根；
4. 从根到叶子的路径上，左分支记 \texttt{0}、右分支记 \texttt{1}，得到每个符号的码字。

\textbf{性质.}

\begin{itemize}
\item Huffman 码是\textbf{前缀码}；
\item 在所有前缀码中，Huffman 码的\textbf{平均码长最小}；
\item 最优平均码长满足
\end{itemize}

\[
  H(P)\le L^*<H(P)+1.
\]

\textbf{最优性证明思路（交换论证）.} 在最优前缀码树中，概率最小的两个符号一定可以放在最深层的兄弟位置：如果它们不在最深处，把它们与最深处的叶子交换，平均码长不会增加。于是可以把这两个最小概率合并成一个“超符号”，递归构造；这正是 Huffman 算法。由此归纳可得最优性。

\textbf{例子（思路）.} 若概率为

\[
0.4,\ 0.2,\ 0.2,\ 0.1,\ 0.1,
\]

则先合并两个 $0.1$ 得到 $0.2$，再合并两个 $0.2$ 得到 $0.4$，最后合并 $0.2$ 与 $0.4$；一种最优码长是 $2,2,2,3,3$（不同合并顺序可能得到不同但等价的树）。概率大的符号码短，概率小的码长。

\subsection{变长码的极限}


对 $n$ 个独立同分布符号一起做前缀编码，最优平均码长 $L_n^*$ 满足

\[
nH(P)\le L_n^*<nH(P)+1.
\]

因此每符号平均码长

\[
\frac{L_n^*}{n}\to H(P).
\]

这就是“熵是可逆压缩的极限速率”。

\subsection{常见错误}


\begin{itemize}
\item 把 Kraft 不等式写成“$\sum2^{-l_i}=1$”；它是不等式，$\le1$ 即可。
\item 认为 Huffman 对\textbf{任意}唯一可译码最优；它是在\textbf{前缀码}类中最优。前缀码类已经能达到唯一可译码的最优渐近速率。
\item 把“最短码”给最不常出现的符号；恰好相反，最短码给概率最大的符号。
\item 忘记 Huffman 需要已知分布；分布未知时要用通用压缩/随机压缩（下一节）。
\end{itemize}
